\chapter{Term Graph Rewriting, Sharing, and Cycles}
\label{ch:term-graph-rewriting}

The tree $\operatorname{pair}(u,u)$ has two occurrences of u.
A graph may represent both by one node with two incoming edges.
Updating that node changes both unfolded occurrences at once.
This is the mechanism behind useful sharing, but it is also the reason a graph rewrite is not merely a faster implementation of an arbitrary tree rewrite.

\section{Define the graph and its read-back}
\label{sec:ch86-graphs}
A rooted term graph has labeled nodes whose ordered outgoing edges match symbol arities.
Several edges may point to the same node.
For an acyclic graph, \term{unfolding} recursively follows edges to produce a tree, duplicating a shared subtree whenever it is reached through different incoming paths.
The unfolding can be exponentially larger than the graph.

For a chain $n_{i+1}=\operatorname{pair}(n_i,n_i)$ above a leaf $n_0$, the graph has $h+1$ nodes at height h.
Its unfolded tree has $2^{h+1}-1$ nodes.
Table~\ref{tab:ch86-growth} works the first heights.
A graph algorithm bounded by graph size must not quietly materialize the whole unfolding inside its claimed constant-time equality operation.

\begin{table}[htbp]
\caption{Sharing can compress exponentially many tree occurrences into a short graph.}
\label{tab:ch86-growth}
\centering
\begin{tabular}{rrr}
\toprule
Height & Graph nodes & Unfolded tree nodes \\
\midrule
0 & 1 & 1 \\
1 & 2 & 3 \\
2 & 3 & 7 \\
3 & 4 & 15 \\
\bottomrule
\end{tabular}
\end{table}

Algorithm~\ref{alg:ch86-unfold} performs an explicit finite read-back and rejects a cycle on the active recursion path.
Its time and output storage are proportional to unfolded output size, not just the number of distinct heap nodes.
The active-path set detects cycles without mistaking harmless repeated visits through separate paths for cycles.
This is the read-back implemented in the lab; it deliberately does not interpret cyclic graphs as infinite rational terms.

\begin{algorithm}[htbp]
\caption{Read back a finite acyclic term graph, rejecting cyclic tree unfolding.}
\label{alg:ch86-unfold}
\begin{pseudocode}
\Require valid node IDs, ordered arity-correct children, root
\Ensure unfolded tree or cycle diagnostic
\Function{Unfold}{$n,active$}
\If{$n\in active$} \State reject cyclic unfolding \EndIf
\State add $n$ to the path-local active set
\If{$n$ is an indirection} \State return unfolding of its target \EndIf
\State recursively unfold children with this active set
\State return labeled tree node containing those child trees
\EndFunction
\end{pseudocode}
\end{algorithm}

\section{Distinguish node update from occurrence replacement}
\label{sec:ch86-update}
The lab graph has nodes v, zero, add(v,zero), and pair(add,add).
One add node is shared by both pair edges.
Applying add-zero in place replaces the add node with an indirection to v.
Unfolding afterward yields pair(v,v); both occurrences changed.
Figure~\ref{fig:ch86-sharing} shows the topology before the update.

\begin{figure}[htbp]
\centering
\begin{tikzpicture}[node distance=13mm and 15mm]
\node[basic block] (p) {pair};\node[basic block,below=of p] (a) {shared add};
\node[basic block,below left=of a] (v) {v};\node[basic block,below right=of a] (z) {zero};
\draw[cfg edge] (p.south west) to[bend right=20] (a.west);
\draw[cfg edge] (p.south east) to[bend left=20] (a.east);
\draw[cfg edge] (a)--(v);\draw[cfg edge] (a)--(z);
\end{tikzpicture}
\caption{Two parent edges identify one shared redex, so a node update affects both unfolded occurrences.}
\label{fig:ch86-sharing}
\end{figure}

If the intended semantics changes only the first occurrence, use path copying or split sharing before the update.
If the intended semantics is call-by-need, updating a shared thunk is deliberate: later uses reuse the same computed value.
For an effectful action, merging two originally separate computations can reduce the number of effects and change behavior.
Sharing therefore requires a semantic justification, not merely a structural hash collision check.

Algorithm~\ref{alg:ch86-update} states the executed pure identity update.
The graph must be a valid acyclic arithmetic term graph, the selected node must have add arity two, and the second child must denote the zero constant.
Only under the pure total identity contract is replacing the node with an indirection sound.
The actual local update costs $O(1)$; maintaining reverse uses, analysis facts, and reclaiming now-dead nodes adds separate work.

\begin{algorithm}[htbp]
\caption{Redirect a shared pure add-zero node while preserving all incoming edges.}
\label{alg:ch86-update}
\begin{pseudocode}
\Require valid acyclic pure arithmetic graph; add node $n$ with children $(a,zero)$
\Ensure every use of $n$ denotes $a$ under the identity contract
\State replace node $n$'s payload with an indirection to $a$
\State invalidate facts that depend on $n$'s former operator or children
\State account for removed child references under the heap's ownership policy
\end{pseudocode}
\end{algorithm}

The small lab implements payload redirection and tests read-back; it does not implement a collector or reverse-use maintenance.
It retains all allocated nodes in its finite test heap, so unreclaimed zero nodes are an explicit model boundary, not a claim of leak-free production graph reduction.
The runtime chapter develops the actual ownership and reclamation obligations for a language implementation.

\section{Preserve meaning under nonlinearity and matching}
\label{sec:ch86-nonlinear}
A tree pattern pair(x,x) asks whether the two subject subterms are equal under the pattern's equality notion.
A graph matcher that tests pointer identity recognizes existing sharing, not all equal unfoldings.
Two distinct nodes can represent equal finite trees; one node can also participate in a cyclic graph whose read-back is outside the finite-term domain.
Choosing identity, structural equality, bisimulation, or equational equality changes the matching problem.

Table~\ref{tab:ch86-equalities} distinguishes these relations.
Hash-consing can make finite structural equality coincide with node identity for canonical immutable acyclic nodes, but only if every construction passes through the interning discipline.
Destructive updates can invalidate a hash-cons key, so mutable reducer heaps and immutable hash-consed DAGs should not share one undocumented invariant.
Indirections need normalization or explicit treatment in both equality and indexing.

\begin{table}[htbp]
\caption{Graph identity and represented-term equality answer different questions.}
\label{tab:ch86-equalities}
\centering
\begin{tabular}{ll}
\toprule
Relation & Meaning \\
\midrule
node identity & exactly the same allocated node \\
finite structural equality & equal acyclic tree unfoldings \\
bisimulation & matching labeled behavior including allowed cycles \\
rewrite convertibility & connected by the chosen equational theory \\
observational equivalence & indistinguishable in permitted program contexts \\
\bottomrule
\end{tabular}
\end{table}

Graph/term correspondence theorems require restrictions on rules, graph interpretation, and reduction strategy.
For example, sharing a variable occurrence on the right side preserves one computation rather than copying its operational work.
The term-graph literature studies precisely when that implementation simulates term rewriting; those results must not be generalized to arbitrary effectful graphs or arbitrary node merging~\cite{barendregt-1987-term-graphs}.

\section{Give cycles an explicit operational meaning}
\label{sec:ch86-cycles}
A cycle can encode recursion or an infinite data structure, but a finite tree reader cannot simply recurse through it.
A lazy list such as ones can be represented by a constructor whose tail points back to the same evaluated list.
Forcing the outer constructor can terminate even though fully unfolding the list cannot.
A self-dependent unevaluated thunk, by contrast, can reenter its own evaluation and require black-hole detection.

Figure~\ref{fig:ch86-cycle} contrasts these roles.
The distinction is whether observation reaches a constructor before recurring, not whether the heap contains any cycle at all.
A collector must also reclaim unreachable cycles; plain reference counting without a cycle strategy can retain them forever.
Recursive closure environments in the functional lab rely on Python's managed heap rather than claiming a new collector implementation.

\begin{figure}[htbp]
\centering
\begin{tikzpicture}[node distance=28mm]
\node[basic block] (cons) {evaluated Cons};
\node[basic block,right=of cons] (thunk) {evaluating thunk};
\draw[cfg edge] (cons) edge[loop below] node[below,font=\small]{tail points back}(cons);
\draw[cfg edge,dashed] (thunk) edge[loop below] node[below,font=\small]{reenters before value}(thunk);
\end{tikzpicture}
\caption{A productive constructor cycle and an evaluation black hole require different treatment.}
\label{fig:ch86-cycle}
\end{figure}

\section{Separate ordinary sharing from optimal reduction}
\label{sec:ch86-optimal}
Sharing one thunk across its uses avoids recomputing that particular suspended expression.
It does not automatically implement optimal lambda reduction, which concerns avoiding certain duplicated reduction work across families of residuals.
Sharing graphs with explicit duplication management introduce additional node kinds and correctness invariants.
Their graph bookkeeping can itself be expensive; optimality in a formal reduction metric is not a universal wall-clock performance guarantee.

Asperti and Guerrini provide background on that larger design space~\cite{asperti-1998-optimal}, while the executable lab intentionally implements only finite read-back and one ordinary shared update.
The functional evaluator implements call-by-need memoization, not a Lamping-style optimal reducer.
Keeping those claims separate lets the reader extend a small correct model without inheriting an unsupported optimality label.

\section{Track residual work after duplication and erasure}
\label{sec:ch86-residuals}
A redex can move, duplicate, or disappear when an enclosing redex is contracted.
Its surviving counterparts are often described as residuals.
For f(x)$\to$pair(x,x), reducing the outer f duplicates the occurrences of a tree redex inside the substituted x.
For f(x)$\to$c, the same inner work is erased.
An evaluator that always reduces arguments before the outer call can therefore perform work that an outermost or lazy strategy never needs.

A shared graph can preserve one underlying x node across both pair edges, so updating that node realizes the effect of reducing both corresponding tree residuals.
That explains the performance opportunity and the simulation complication at the same time.
Counting one graph update as exactly one tree rewrite is generally wrong when there are multiple incoming paths.
Correctness should relate represented values or an appropriately defined reduction simulation, while performance should state which operations are being counted.

The distinction also appears in garbage collection.
An erasing rewrite removes references to an argument subgraph but does not prove that the argument is globally unreachable.
Other graph roots or shared parents may retain it.
A reference-count implementation decrements the removed edges and may trigger recursive release; a tracing collector discovers reachability from the complete root set.
Neither should free a node merely because one redex stopped using it.

\section{Maintain indirections and sharing without corrupting indexes}
\label{sec:ch86-indirections}
Indirections let an update preserve node identity while changing the value it denotes.
Chains of indirections can accumulate after successive rewrites.
Path compression can shorten future traversals, but a reducer must respect the heap's mutation, GC, and concurrency rules when rewriting those links.
A self-indirection is not a valid finite acyclic shortcut and is rejected by the lab's read-back.
The shape of indirections should be part of heap verification rather than an assumption buried in the evaluator.

Hash-consing poses a different issue.
An intern table keyed by an immutable operator and child tuple is valid only while that tuple remains immutable.
Updating an interned node in place without removing or repairing its old table entry can make later construction retrieve a node with the wrong structure.
One design separates immutable syntax DAGs from mutable evaluation nodes; another uses stable handles to replace immutable payloads while maintaining explicit versioned indexes.
Both require a deliberate invariant, not merely a shared pointer type.

Finally, graph matching must say whether it preserves or reflects sharing.
A homomorphic match may map distinct pattern nodes to one host node; an injective match forbids that identification.
Those choices affect deletion and duplicated-variable behavior.
The shared arithmetic fixture performs neither general search nor arbitrary graph homomorphism: the caller supplies one node in a known finite graph.
That narrow contract makes the executed update easy to audit while leaving the broader matching design visible.

\section{Exercises}
\label{sec:ch86-exercises}
\begin{enumerate}
\item\label{ex:ch86-unfold} Compute graph and unfolded sizes at height four in the doubled-pair family.
\item\label{ex:ch86-update} Unfold the lab's pair/shared-add graph before and after redirecting add-zero.
How many tree occurrences change?
\item\label{ex:ch86-identity} Two distinct immutable nodes both represent f(a).
Explain why pointer inequality does not imply unequal tree patterns.
\item\label{ex:ch86-cycle} Contrast a constructor tail cycle with a thunk that immediately forces itself.
Which can expose weak-head structure, and why does finite read-back reject both cycles?
\end{enumerate}

\section{Further reading}
\label{sec:ch86-reading}
Barendregt and colleagues study the graph/term correspondence under explicit restrictions~\cite{barendregt-1987-term-graphs}; Ohlebusch provides the broader rewrite-theoretic context~\cite{ohlebusch-2002-rewriting}.
The STG paper shows how sharing and updates become runtime mechanisms for non-strict functional execution~\cite{peytonjones-1992-stg}.
